\documentclass[11pt]{article}
\usepackage[margin=1in]{geometry}
\usepackage{amsmath, amssymb, amsfonts}
\usepackage{bm}
\usepackage{graphicx}
\usepackage{hyperref}
\usepackage{enumitem}
\usepackage{booktabs}

\title{Plato's Cave and The Genius:\\
Collective Ascent, Trust Dynamics, and Meta-Reform \\ Beyond Self-Correcting Education}
\author{Kara Olivarria et al. \\ \\ A Citizen Gardens Research Initiative \\ The Foundation of Multiplicity}
\date{April 8, 2026}

\begin{document}
\maketitle

\begin{abstract}
This report presents a computational realization of Plato's Cave as a
multi-agent sandbox for simulating the emergence of hierarchical trust
and the collective identification of latent truths, grounded in
Kara Olivarria's 25-dimensional quantum education architecture.
The sandbox couples a latent \emph{Form}---a 25D prime-indexed
pedagogical lattice---with agents that observe noisy shadows of the Form
and update their beliefs via local sparse refinement and a social
trust--oracle mechanism.
On top of this, the ``Genius'' model is used as a meta-architecture:
its five-phase cycle (novelty $\rightarrow$ enhancement $\rightarrow$
critique $\rightarrow$ formalization $\rightarrow$ validation) governs
structural meta-reforms, including hierarchy evolution and
dimension-level mutations.
We formalize the underlying mathematics, specify implementable
algorithms, and report synthetic validation experiments demonstrating
epistemic ascent, suppression of ``false suns,'' and earned hierarchical
authority.
\end{abstract}
\newpage
\tableofcontents
\newpage
\section{Executive Summary}

We construct a computational interpretation of Plato's Cave in which
agents attempt to recover an invariant latent \emph{Form} from partial,
noisy projections (shadows).
The latent Form is instantiated as Olivarria's 25-dimensional quantum
education architecture, a prime-indexed pedagogical lattice
representing a planetary curriculum.

Each agent $i$ holds a state $\bm{x}^{(i)}(t) \in \mathbb{R}^{25}$ and
observes shadows
\[
  \bm{y}_i = \bm{A}_i \bm{\theta}^* + \bm{\varepsilon}_i,
\]
where $\bm{\theta}^* \in \mathbb{R}^{25}$ is the fixed Form and
$\bm{A}_i$ encodes the agent's limited view.
Agents perform local LASSO-like refinement and then
enter a social dialectic where trust is assigned via
counterfactual explanatory power and regularized by an adaptive oracle
ensemble.

On a slower timescale, a ``Genius'' meta-architecture governs
hierarchy evolution and structural mutations.
Agents earn higher tiers through an impact--atypicality index, and
dimensions are added, split, or pruned when Genius-driven metrics
over dimension usage and oracle contributions cross thresholds.

Synthetic experiments show that, even with modest noise and adversarial
agents, the system:
\begin{itemize}[noitemsep]
  \item reduces mean distance to the 25D Form over time;
  \item suppresses ``false suns'' (locally coherent but globally
        misaligned subspaces);
  \item yields stable, earned hierarchies aligned with oracle performance;
  \item supports principled architecture evolution under the Genius cycle.
\end{itemize}

\section{The 25D Form: Quantum Education Architecture}

\subsection{Latent Form and Prime-Indexed Structure}

We model Olivarria's 25-dimensional quantum education architecture as
a latent vector
\[
  \bm{\theta}^* \in \mathbb{R}^{25},
\]
where each coordinate corresponds to a macro-dimension of the
curriculum (e.g., Primordial Cognitive Lattice, Apophatic Learning
Singularity, Quantum Institutional DNA, etc.).
At a finer level, each coordinate can be understood as a
prime-indexed multiplicity structure, but for computational purposes
we treat the macro-form as a 25D vector.

Optionally, we can represent an agent's curriculum state using a
multiplicative prime code:
\[
  P^{(i)}(t) = \prod_{k=1}^{25} p_k^{x^{(i)}_k(t)},
\]
where $p_k$ is the $k$-th prime.
In practice, all computation is carried out in exponent space
$\bm{x}^{(i)}$, and the prime-product encoding is reserved for
interpretation and logging.

\subsection{Observation Model: Shadows in the Cave}

Agent $i$ does not observe $\bm{\theta}^*$ directly.
Instead, it observes a shadow
\[
  \bm{y}_i = \bm{A}_i \bm{\theta}^* + \bm{\varepsilon}_i,
\]
where $\bm{A}_i \in \mathbb{R}^{m_i \times 25}$ encodes a partial,
possibly biased projection, and $\bm{\varepsilon}_i$ is noise.

When comparing candidate states $\bm{x} \in \mathbb{R}^{25}$ against
agent $i$'s observations, we define a local loss
\[
  E_i(\bm{x}) = \frac{1}{m_i} \left\| \bm{A}_i \bm{x} - \bm{y}_i \right\|_2^2
               + \lambda \| \bm{x} \|_1,
\]
where the $\ell_1$ term encourages sparsity in the 25D representation
relative to the Form.

\section{Local Learning: LASSO Refinement}

\subsection{Proximal Gradient Dynamics}

Each agent updates its current state $\bm{x}^{(i)}(t)$ by a small number
of proximal gradient steps on $E_i$:
\begin{align*}
  \bm{g}_i(t) &= \frac{1}{m_i} \bm{A}_i^\top
                \bigl( \bm{A}_i \bm{x}^{(i)}(t) - \bm{y}_i \bigr), \\
  \bm{z}_i(t+1) &= \bm{x}^{(i)}(t) - \eta \bm{g}_i(t), \\
  \bm{x}^{(i)}(t+1) &= \operatorname{prox}_{\eta \lambda \|\cdot\|_1}
                      \bigl( \bm{z}_i(t+1) \bigr),
\end{align*}
where $\eta > 0$ is a learning rate and the proximal operator is
soft-thresholding:
\[
  \operatorname{prox}_{\tau \|\cdot\|_1}(\bm{z})
  = \operatorname{sign}(\bm{z}) \odot
    \max\bigl( |\bm{z}| - \tau, \, 0 \bigr).
\]

In the full system, agents typically run a few such steps per outer
iteration before entering the social dialectic.

\section{Social Dialectic: Trust and Oracle Ensemble}

\subsection{Counterfactual Loss and Normalization}

Let $\tilde{\bm{x}}^{(j)}$ denote agent $j$'s current refined state
after local updates.
Agent $i$ evaluates how well $j$'s state explains its own shadow via
the counterfactual loss
\[
  L_{ij}(t) = E_i\bigl( \tilde{\bm{x}}^{(j)}(t) \bigr).
\]
To avoid scale distortions across agents with different shadow
variances, we normalize:
\[
  L_{ij}^{\text{norm}}(t) =
  \frac{L_{ij}(t)}{\widehat{\mathrm{Var}}(\bm{y}_i) + \varepsilon},
\]
where $\widehat{\mathrm{Var}}(\bm{y}_i)$ is an empirical variance estimate
and $\varepsilon > 0$ is a small constant.

\subsection{Softmax Trust with Annealing and Inertia}

We interpret $-L_{ij}^{\text{norm}}$ as an unnormalized log-likelihood
of $j$'s model from $i$'s perspective.
For each agent $i$ we define a trust distribution over senders $j$:
\[
  S_{ij}^{\mathrm{raw}}(t)
  = \frac{
      \exp\!\bigl(-\alpha_t L_{ij}^{\text{norm}}(t)\bigr)
    }{
      \sum_{k} \exp\!\bigl(-\alpha_t L_{ik}^{\text{norm}}(t)\bigr)
    },
\]
where $\alpha_t$ is an annealed sensitivity parameter that increases
over time.
To model epistemic inertia, we use an exponential moving average
(EMA) over trust matrices:
\[
  S_{ij}(t) = \gamma S_{ij}(t-1)
             + (1-\gamma) S_{ij}^{\mathrm{raw}}(t),
\]
with $\gamma \in (0,1)$.

This softmax–EMA mechanism approximates Bayesian pooling; calibration
can be monitored via proper scoring rules such as the Brier score.

\subsection{Oracle Ensemble and Invariance Penalty}

We introduce a set of oracle functions $\{ O_o \}_{o=1}^K$ that
evaluate candidate states $\bm{x}$ against separate criteria (e.g.,
overall fit to held-out data, equity measures, long-term retention).
For agent $j$, the oracle costs at time $t$ are:
\[
  c_{j,o}(t) = O_o\bigl( \tilde{\bm{x}}^{(j)}(t) \bigr).
\]

We maintain an exponentially weighted history of oracle costs and define
a per-oracle degradation
\[
  \Delta_{j,o}(t) = c_{j,o}(t) - \bar{c}_{j,o}(t-1),
\]
where $\bar{c}_{j,o}$ is a recent average.
The oracle penalty for agent $j$ is then
\[
  \pi_j(t) = \exp\!\left(
    -\kappa
    \sum_{o=1}^K \phi_o(t)
    \frac{\max(\Delta_{j,o}(t), 0)}{\tau_\pi}
  \right),
\]
with $\phi_o(t)$ an adaptive weight reflecting oracle $o$'s predictive
usefulness, $\kappa > 0$, and $\tau_\pi > 0$ a scale.
We enforce a floor $\pi_j(t) \geq \pi_{\min}$.

Trust utilities are modified by oracle penalties, e.g.
\[
  U_{ij}(t) = - L_{ij}^{\text{norm}}(t) + \beta_h h^{(j)}(t),
\]
then
\[
  S_{ij}^{\mathrm{raw}}(t)
  \propto \exp\bigl(\alpha_t U_{ij}(t)\bigr) \, \pi_j(t).
\]

\subsection{Social Mixing}

Given the row-stochastic trust matrix $S(t)$, we update states via
social mixing:
\[
  \bm{x}^{(i)}_{\text{social}}(t+1)
  = \sum_{j} S_{ij}(t) \, \tilde{\bm{x}}^{(j)}(t).
\]
We may blend local and social updates:
\[
  \bm{x}^{(i)}(t+1)
  = (1-\beta_{\text{mix}}) \tilde{\bm{x}}^{(i)}(t)
    + \beta_{\text{mix}} \bm{x}^{(i)}_{\text{social}}(t+1),
\]
with $\beta_{\text{mix}} \in [0,1]$.

\section{Hierarchy Evolution and the Genius Index}

\subsection{Hierarchy Beliefs}

Each agent $i$ maintains a hierarchy belief
$h^{(i)}(t) \in [0,1]$ representing its institutional weight.
This is used as a bias in trust formation and oracle subsampling.

We define:
\begin{itemize}[noitemsep]
  \item Oracle alignment:
    \[
      o_i(t) = \exp(-\tau_o \pi_i(t)),
    \]
  \item Weighted centrality:
    \[
      c_i(t) = \sum_j S_{ji}(t) h^{(j)}(t),
      \quad
      c_i^{\text{norm}}(t) = \frac{c_i(t)}{\max_k c_k(t) + \varepsilon}.
    \]
\end{itemize}

\subsection{Genius Index}

The Genius model views a ``genius'' as a trajectory whose moves are both
high-impact and atypical relative to a baseline distribution.
We approximate this at the agent level by defining
\[
  G_i(t) = (1-\lambda) o_i(t) + \lambda c_i^{\text{norm}}(t),
\]
and an atypicality measure $A_i(t)$ derived from the agent's prime-move
usage (or, in the minimal implementation, a proxy based on variance or
surprise).
The Genius index is
\[
  J_i(t) = G_i(t) \, \bigl( 1 + \eta A_i(t) \bigr),
\]
with $\eta > 0$ small.

\subsection{Hierarchy Update}

We map $J_i(t)$ into a target hierarchy level via a sigmoid:
\[
  h_{\text{target}}^{(i)}(t)
  = \sigma\bigl( \beta_J J_i(t) \bigr)
  = \frac{1}{1 + e^{-\beta_J J_i(t)}},
\]
and update $h^{(i)}$ sluggishly:
\[
  h^{(i)}(t+1)
  = \eta_h h_{\text{target}}^{(i)}(t)
    + (1-\eta_h) h^{(i)}(t),
\]
where $\eta_h \in (0,1)$ controls hierarchy adaptation speed.

Agents with $h^{(i)}(t)$ above a promotion threshold $\theta_{\text{promote}}$
are treated as higher-tier nodes in subsequent iterations, affecting
trust bias and oracle sampling.

\section{Genius Meta-Architecture and Meta-Reform}

\subsection{The Five-Phase Genius Cycle}

The Genius model provides a meta-architecture with phases:
\begin{enumerate}[noitemsep]
  \item \textbf{Novelty}: propose new moves (e.g., architectural mutations).
  \item \textbf{Enhancement}: locally refine and implement them.
  \item \textbf{Critique}: stress-test via oracle and social dynamics.
  \item \textbf{Formalization}: freeze successful moves into structure.
  \item \textbf{Validation}: benchmark against baselines and prior phases.
\end{enumerate}

We apply this cycle at the system level to govern structural changes to
the 25D architecture and the trust/hierarchy mechanisms themselves.

\subsection{Dimension-Level Mutations (Outline)}

Briefly, for each dimension $k$ we can track:
\begin{itemize}[noitemsep]
  \item Usage and contribution to oracle fit;
  \item Genius-weighted innovations involving that dimension;
  \item Redundancy with other dimensions.
\end{itemize}
A dimension with low contribution and low Genius-linked activity may be
pruned; overloaded dimensions may be split into subdimensions; new
dimensions may be proposed by high-Genius agents.
Each such mutation is treated as a novelty proposal, run through
enhancement/critique/formalization/validation, and either adopted or
rolled back.

\section{Minimal Algorithm and Code Snippet}

\subsection{High-Level Algorithm}

At a high level, one outer iteration of the sandbox proceeds as:
\begin{enumerate}[noitemsep]
  \item Local LASSO refinement for each agent (fast inner loop).
  \item Compute counterfactual losses and normalized trust.
  \item Apply oracle ensemble penalties to adjust trust utilities.
  \item Update trust matrix via softmax + EMA.
  \item Socially mix agent states via trust.
  \item Update hierarchy beliefs via Genius index.
  \item (Occasionally) propose and evaluate meta-reforms under the
        Genius cycle.
\end{enumerate}

\subsection{Compact NumPy-Style Snippet}

Below is a compact Python/NumPy-style snippet capturing the core trust
and mixing dynamics (for a fixed architecture):

\begin{verbatim}
import numpy as np

def proximal_lasso_step(x, A_i, y_i, lr, lam):
    # x: (d,), A_i: (m, d), y_i: (m,)
    residual = A_i @ x - y_i
    grad = A_i.T @ residual / len(y_i)
    z = x - lr * grad
    # soft-thresholding
    return np.sign(z) * np.maximum(np.abs(z) - lr * lam, 0.0)

def trust_oracle_update(X, A, Y, alpha_t, gamma,
                        oracle_penalty, oracle_penalty_tau,
                        penalty_min, rng):
    N, d = X.shape
    # counterfactual losses
    # pred[i,j,:] = A[i] @ X[j]
    pred = np.einsum('imd,jd->ijm', A, X)
    losses = np.mean((pred - Y[:, None, :])**2, axis=2)
    var_y = np.var(Y, axis=1, keepdims=True) + 1e-8
    L_norm = losses / var_y

    row_min = np.min(L_norm, axis=1, keepdims=True)
    base = np.exp(-alpha_t * (L_norm - row_min))

    # simple oracle: global subsample
    k = max(5, int(0.2 * N))
    idx_sub = rng.choice(N, k, replace=False)
    pred_sub = np.einsum('kmd,jd->kjm', A[idx_sub], X)
    err_sub = pred_sub - Y[idx_sub, None, :]
    E_oracle = np.mean(err_sub**2, axis=(0, 2))  # (N,)

    delta = E_oracle - np.median(E_oracle)
    penalty = np.exp(-oracle_penalty
                     * np.maximum(delta, 0.0)
                     / oracle_penalty_tau)
    penalty = np.maximum(penalty, penalty_min)
    score = base * penalty[None, :]

    # row-stochastic trust
    score = score - np.max(score, axis=1, keepdims=True)
    S = np.exp(score)
    S = S / (np.sum(S, axis=1, keepdims=True) + 1e-12)
    S = gamma * S + (1 - gamma) * S  # EMA placeholder
    X_new = S @ X
    return X_new, S, E_oracle
\end{verbatim}

This snippet can be embedded in a full simulation loop that also
performs local LASSO updates and hierarchy evolution.

\section{Synthetic Experiments (Outline)}

\subsection{Setup}

We consider $N$ agents, each with an observation matrix $\bm{A}_i$ and
noise level $\sigma$.
We initialize $\bm{\theta}^*$ as a sparse 25D vector and generate
shadows $\bm{y}_i$.
We compare:
\begin{itemize}[noitemsep]
  \item local-only learning,
  \item trust-only social mixing,
  \item combined local+trust+oracle+hierarchy dynamics.
\end{itemize}

We also seed a fraction of agents as ``false suns'' with initial states
aligned to a wrong subspace and track their trust centrality and
hierarchy beliefs.

\subsection{Key Metrics}

\begin{itemize}[noitemsep]
  \item Mean distance to the Form:
    \(
      \frac{1}{N}\sum_i \| \bm{x}^{(i)}(t) - \bm{\theta}^* \|.
    \)
  \item Mean oracle penalty.
  \item Trust centrality of false suns:
    sum of incoming trust weights from non-adversarial agents.
  \item Hierarchy entropy over time.
\end{itemize}

Observed trends in preliminary runs include steady descent in mean
distance to the Form, decreasing trust centrality for false suns, and
stabilization of hierarchy entropy at an ordered level.

\section{Discussion and Future Work}

The Plato--25D--Genius sandbox operationalizes a philosophical allegory
as a testable, multi-layered dynamical system:
local sparse refinement, social trust and oracle-based correction,
and Genius-driven meta-reform.
Future work includes:
\begin{itemize}[noitemsep]
  \item richer oracle ensembles including real-world signals;
  \item explicit modeling of the Form as a manifold
        $\mathcal{M} \subset \mathbb{R}^{25}$;
  \item phase-based architecture evolution with formal convergence and
        robustness guarantees;
  \item empirical mapping to real educational data.
\end{itemize}

%%%%%%%%%%%%%%%%%%%%%%%%%%%%%%%%%%%%%%%%%%%%%%%%%%%%%%%%%%%%
% MATHEMATICAL APPENDIX
%%%%%%%%%%%%%%%%%%%%%%%%%%%%%%%%%%%%%%%%%%%%%%%%%%%%%%%%%%%%

\appendix
\section{Mathematical Appendix}

In this appendix we formalize several of the claims made in the main
text and provide explicit operator-norm bounds under standard
assumptions.
We focus on:
(i) stability and descent of the local LASSO refinement,
(ii) contraction properties of the trust-based social mixing,
and (iii) interaction of the trust–oracle mechanism with false-sun
clusters under bounded noise and restricted operator norms.

Throughout, we adopt the following notation:
$\|\cdot\|$ denotes the Euclidean norm on $\mathbb{R}^d$,
$\|\cdot\|_{\mathrm{op}}$ the induced operator norm on matrices,
and $\bm{1}$ the all-ones vector of appropriate dimension.

\subsection{Assumptions}

We work under the following assumptions.

\begin{enumerate}[label=(A\arabic*),noitemsep]
  \item \textbf{Form and observations.}
    The latent Form is fixed and given by
    $\bm{\theta}^* \in \mathbb{R}^{d}$ with $d=25$.
    For each agent $i \in \{1,\dots,N\}$ there is an observation matrix
    $\bm{A}_i \in \mathbb{R}^{m_i \times d}$ and noisy observations
    \[
      \bm{y}_i = \bm{A}_i \bm{\theta}^* + \bm{\varepsilon}_i,
    \]
    where $\bm{\varepsilon}_i$ are independent, mean-zero random
    vectors with
    $\mathbb{E}\|\bm{\varepsilon}_i\|^2 \leq m_i \sigma^2$ for some
    $\sigma > 0$.

  \item \textbf{Operator-norm bounds.}
    There exist constants $0 < \mu_{\min} \leq \mu_{\max} < \infty$
    such that for all $i$,
    \[
      \mu_{\min}
      \;\leq\;
      \|\bm{A}_i\|_{\mathrm{op}}^2
      \;\leq\;
      \mu_{\max}.
    \]

  \item \textbf{Step size for local refinement.}
    The learning rate $\eta > 0$ for the local gradient step satisfies
    \[
      0 < \eta < \frac{2}{\mu_{\max}}.
    \]

  \item \textbf{Trust matrix structure.}
    At each outer iteration $t$, the trust matrix
    $\bm{S}(t) \in \mathbb{R}^{N \times N}$ is row-stochastic:
    \[
      S_{ij}(t) \geq 0,
      \qquad
      \sum_{j=1}^N S_{ij}(t) = 1 \quad \forall i.
    \]
    We denote by
    $\rho\bigl(\bm{S}(t)\bigr)$ the spectral radius of $\bm{S}(t)$
    considered as a linear operator on the subspace orthogonal to
    $\bm{1}$.

  \item \textbf{Oracle alignment and penalty floor.}
    For each $t$, the oracle penalties $\pi_j(t)$ satisfy
    $\pi_j(t) \in [\pi_{\min},1]$ with $\pi_{\min} \in (0,1)$ fixed.
    The modified utilities used to construct $\bm{S}(t)$ are such that
    for non-adversarial agents, expected oracle losses decrease in $t$.
\end{enumerate}

\subsection{Local LASSO Refinement: Stability and Descent}

We consider the local objective for agent $i$:
\[
  E_i(\bm{x})
  =
  \frac{1}{m_i}
  \left\|
    \bm{A}_i \bm{x} - \bm{y}_i
  \right\|_2^2
  + \lambda \|\bm{x}\|_1,
\]
with $\lambda \geq 0$.
The proximal gradient update for agent $i$ at time $t$ is
\begin{align*}
  \bm{g}_i(t)
    &= \frac{1}{m_i} \bm{A}_i^\top
       \bigl( \bm{A}_i \bm{x}^{(i)}(t) - \bm{y}_i \bigr),
\\
  \bm{z}_i(t+1)
    &= \bm{x}^{(i)}(t) - \eta \bm{g}_i(t),
\\
  \bm{x}^{(i)}(t+1)
    &= \operatorname{prox}_{\eta \lambda \|\cdot\|_1}
       \bigl( \bm{z}_i(t+1) \bigr).
\end{align*}

\begin{lemma}[Lipschitz constant of the smooth part]
Under (A2), the gradient of the smooth part
\[
  f_i(\bm{x})
  = \frac{1}{m_i}\bigl\|\bm{A}_i \bm{x} - \bm{y}_i\bigr\|_2^2
\]
is Lipschitz with constant
$L_i = \frac{2}{m_i} \|\bm{A}_i\|_{\mathrm{op}}^2$.
In particular, $L_i \le \frac{2}{m_i} \mu_{\max}$.
\end{lemma}

\begin{proof}
We have
\[
  \nabla f_i(\bm{x})
  = \frac{2}{m_i} \bm{A}_i^\top
    \bigl( \bm{A}_i \bm{x} - \bm{y}_i \bigr).
\]
For any $\bm{x},\bm{z} \in \mathbb{R}^d$,
\begin{align*}
  \|\nabla f_i(\bm{x}) - \nabla f_i(\bm{z})\|
  &=
  \frac{2}{m_i}
  \left\|
    \bm{A}_i^\top \bm{A}_i (\bm{x}-\bm{z})
  \right\|
  \\
  &\le
  \frac{2}{m_i}
  \|\bm{A}_i^\top \bm{A}_i\|_{\mathrm{op}}
  \, \|\bm{x}-\bm{z}\|
  \\
  &=
  \frac{2}{m_i} \|\bm{A}_i\|_{\mathrm{op}}^2
  \, \|\bm{x}-\bm{z}\|.
\end{align*}
Thus, $L_i = \frac{2}{m_i}\|\bm{A}_i\|_{\mathrm{op}}^2$.
\end{proof}

\begin{lemma}[Non-expansiveness of the proximal step]
Let $g(\bm{x}) = \lambda \|\bm{x}\|_1$.
For any $\tau > 0$, the proximal operator
$\operatorname{prox}_{\tau g}$ is firmly non-expansive:
\[
  \bigl\|\operatorname{prox}_{\tau g}(\bm{u})
        -\operatorname{prox}_{\tau g}(\bm{v})\bigr\|
  \le
  \|\bm{u}-\bm{v}\|
  \quad \forall \bm{u},\bm{v}.
\]
\end{lemma}

\begin{proof}
This is a standard property of proximal operators of convex functions.
\end{proof}

\begin{proposition}[Stability of local LASSO refinement]
Under (A2)–(A3), the mapping
\[
  T_i(\bm{x}) =
  \operatorname{prox}_{\eta \lambda \|\cdot\|_1}\!
  \bigl(
    \bm{x} - \eta \nabla f_i(\bm{x})
  \bigr)
\]
is a contraction for each $i$ in a neighborhood of any minimizer of
$E_i$, and the sequence $\bm{x}^{(i)}(t+1) = T_i(\bm{x}^{(i)}(t))$
is stable in the sense that
\[
  \|\bm{x}^{(i)}(t+1)\|
  \leq
  C_i \bigl( 1 + \|\bm{x}^{(i)}(t)\| \bigr),
\]
for some constant $C_i$ depending on $\|\bm{A}_i\|_{\mathrm{op}}$ and
$\|\bm{y}_i\|$.
\end{proposition}

\begin{proof}[Sketch]
Using the Lipschitz continuity of $\nabla f_i$ with constant $L_i$
and the standard proximal-gradient analysis, if
$0 < \eta < 2/L_i$, then $T_i$ is non-expansive and locally contractive
around a minimizer of $E_i$.
The linear growth bound follows from bounding
$\|\nabla f_i(\bm{x})\|$ in terms of $\|\bm{x}\|$ and $\|\bm{y}_i\|$,
and then using the non-expansiveness of the proximal operator.
\end{proof}

\subsection{Trust-Based Social Mixing: Consensus and Norm Bounds}

Let $\bm{X}(t) \in \mathbb{R}^{N \times d}$ be the matrix whose $i$-th
row is $\bm{x}^{(i)}(t)$.
Define the consensus (mean) state
\[
  \bar{\bm{x}}(t)
  = \frac{1}{N}\sum_{i=1}^N \bm{x}^{(i)}(t),
\]
and the centered state matrix
\[
  \bm{E}(t)
  = \bm{X}(t) - \bm{1} \bar{\bm{x}}(t)^\top.
\]
The social mixing update (ignoring local refinement for the moment) is
\[
  \bm{X}(t+1) = \bm{S}(t) \bm{X}(t).
\]

\begin{lemma}[Preservation of the mean]
Under (A4), for any row-stochastic $\bm{S}(t)$,
\[
  \bar{\bm{x}}(t+1) = \bar{\bm{x}}(t).
\]
\end{lemma}

\begin{proof}
\[
  \bar{\bm{x}}(t+1)
  = \frac{1}{N}\sum_{i=1}^N \bm{x}^{(i)}(t+1)
  = \frac{1}{N}\sum_{i=1}^N \sum_{j=1}^N S_{ij}(t)\bm{x}^{(j)}(t)
  = \frac{1}{N}\sum_{j=1}^N \left(
      \sum_{i=1}^N S_{ij}(t)
    \right)\bm{x}^{(j)}(t).
\]
Since $\bm{S}(t)$ is row-stochastic, $\sum_i S_{ij}(t) = 1$ for each $j$,
so
\[
  \bar{\bm{x}}(t+1)
  = \frac{1}{N}\sum_{j=1}^N \bm{x}^{(j)}(t)
  = \bar{\bm{x}}(t).
\]
\end{proof}

\begin{proposition}[Contraction on the disagreement space]
Let $\bm{S}(t)$ satisfy (A4) and assume that, when restricted to the
subspace orthogonal to $\bm{1}$, the operator norm of $\bm{S}(t)$
satisfies
\[
  \bigl\|\bm{S}(t)\bigr\|_{\mathrm{op},\perp \bm{1}}
  \leq \rho < 1
\]
uniformly in $t$.
Then
\[
  \|\bm{E}(t+1)\|_{\mathrm{F}}
  \leq
  \rho \, \|\bm{E}(t)\|_{\mathrm{F}},
\]
where $\|\cdot\|_{\mathrm{F}}$ is the Frobenius norm.
In particular, the disagreement norm decays at least geometrically.
\end{proposition}

\begin{proof}
Observe that
\[
  \bm{E}(t+1)
  = \bm{X}(t+1) - \bm{1}\bar{\bm{x}}(t+1)^\top
  = \bm{S}(t)\bm{X}(t) - \bm{1}\bar{\bm{x}}(t)^\top
  = \bm{S}(t)\bigl(\bm{X}(t) - \bm{1}\bar{\bm{x}}(t)^\top\bigr),
\]
where we used preservation of the mean.
Thus,
\[
  \bm{E}(t+1)
  = \bm{S}(t) \bm{E}(t).
\]
Since each column of $\bm{E}(t)$ lies in the subspace orthogonal to
$\bm{1}$, the operator bound on $\bm{S}(t)$ implies
\[
  \|\bm{E}(t+1)\|_{\mathrm{F}}
  \leq
  \rho \, \|\bm{E}(t)\|_{\mathrm{F}}.
\]
\end{proof}

\subsection{Combined Local and Social Dynamics}

We now consider one outer iteration comprising:
(i) local proximal-gradient refinement for each agent,
(ii) trust/oracle update producing $\bm{S}(t)$,
and (iii) social mixing via $\bm{S}(t)$.

Let $\mathcal{T}$ denote the combined map
\[
  \bm{X}(t) \mapsto \bm{X}(t+1),
\]
obtained by applying local maps $T_i$ row-wise, followed by $\bm{S}(t)$
mixing.
Under (A2)–(A4), for suitably small local step size $\eta$ and
trust matrices satisfying the contraction property above, we obtain a
joint contraction in a weighted norm around a neighborhood of
$\bm{\theta}^* \bm{1}^\top$.
We state a simplified result.

\begin{proposition}[Local contraction around synchronized Form]
Assume (A1)–(A4) and suppose that:
\begin{itemize}[noitemsep]
  \item for each $i$, $T_i$ is a contraction in a neighborhood of
        $\bm{\theta}^*$ with contraction factor $\gamma_i < 1$;
  \item the trust matrices $\bm{S}(t)$ satisfy
        $\|\bm{S}(t)\|_{\mathrm{op},\perp \bm{1}} \leq \rho < 1$
        uniformly.
\end{itemize}
Then there exists a neighborhood $\mathcal{U}$ of the synchronized
state
$\bm{\Theta}^* := \bm{1} (\bm{\theta}^*)^\top$
and a constant $\kappa \in (0,1)$ such that for all
$\bm{X}(t) \in \mathcal{U}$,
\[
  \|\mathcal{T}(\bm{X}(t)) - \bm{\Theta}^*\|_{\mathrm{F}}
  \leq
  \kappa \, \|\bm{X}(t) - \bm{\Theta}^*\|_{\mathrm{F}}.
\]
\end{proposition}

\begin{proof}[Sketch]
We decompose $\bm{X}(t)$ into mean and disagreement parts:
\[
  \bm{X}(t) = \bm{1} \bar{\bm{x}}(t)^\top + \bm{E}(t).
\]
The local maps $T_i$ act independently and, in a neighborhood of
$\bm{\theta}^*$, contract each row toward $\bm{\theta}^*$.
This yields a contraction on the mean component, up to noise.
The trust-based mixing contracts the disagreement component via the
previous proposition.
Combining these two effects yields a joint contraction in a weighted
norm near $\bm{\Theta}^*$.
\end{proof}

\subsection{False-Sun Suppression under Oracle Penalties}

Finally, we formalize the qualitative behavior that ``false suns'' (agents
aligned with an incorrect subspace) are suppressed by the oracle
ensemble.

We model a false-sun agent $j$ as one whose candidate state satisfies
$\|\bm{x}^{(j)} - \bm{\theta}^*\| \geq \delta$ for some fixed
$\delta > 0$ and for all $t$ in a given window (i.e., it remains far
from the true Form).
Assume that the oracle ensemble is aligned with the Form in the sense
that there exists $\gamma > 0$ such that for any agent $i$,
\[
  \mathbb{E}\bigl[ O(\bm{x}^{(i)}(t)) \bigr]
  \geq
  \mathbb{E}\bigl[ O(\bm{\theta}^*) \bigr]
  + \gamma \|\bm{x}^{(i)}(t) - \bm{\theta}^*\|^2,
\]
for a suitable averaged oracle $O$.
Then, for a false-sun agent $j$, we have
\[
  \mathbb{E}\bigl[ O(\bm{x}^{(j)}(t)) \bigr]
  \geq
  \mathbb{E}\bigl[ O(\bm{\theta}^*) \bigr]
  + \gamma \delta^2.
\]

\begin{proposition}[Expected decay of false-sun trust]
Under (A1)–(A5) and the oracle alignment condition above, let
$C_{\mathrm{fs}}(t)$ denote the expected trust centrality of a false-sun
agent $j$ at time $t$, i.e.
\[
  C_{\mathrm{fs}}(t)
  = \mathbb{E}\left[
      \sum_{i=1}^N S_{ij}(t)
    \right].
\]
Then, for sufficiently large $t$ and suitable oracle-penalty parameters,
there exists $\eta_{\mathrm{fs}} \in (0,1)$ such that
\[
  C_{\mathrm{fs}}(t+1)
  \leq
  \eta_{\mathrm{fs}} \, C_{\mathrm{fs}}(t).
\]
\end{proposition}

\begin{proof}[Informal sketch]
The oracle penalty $\pi_j(t)$ for agent $j$ is a decreasing function of
its relative oracle degradation.
For a false-sun agent whose state is persistently far from the Form,
the oracle cost remains strictly worse in expectation than that of
non-adversarial agents.
Thus, $\pi_j(t)$ is bounded away from $1$ and (on average) decreases
relative to other agents.
This enters the utility $U_{ij}(t)$ multiplicatively or additively,
reducing the softmax mass assigned to $j$.
Over time, as the trust matrices are updated, the relative mass in
column $j$ decays geometrically in expectation.
The existence of a uniform factor $\eta_{\mathrm{fs}}<1$ follows from
uniform lower bounds on the gap between false-sun and non-false-sun
oracle costs together with the parameterization of $\pi_j(t)$ in the
softmax.
\end{proof}

This formalizes the qualitative phenomenon observed in synthetic
experiments: ``false suns'' may attract temporary trust if they explain
local shadows well but are ultimately suppressed by the oracle ensemble
as the system converges toward the 25D Form.

%%%%%%%%%%%%%%%%%%%%%%%%%%%%%%%%%%%%%%%%%%%%%%%%%%%%%%%%%%%%
% END MATHEMATICAL APPENDIX
%%%%%%%%%%%%%%%%%%%%%%%%%%%%%%%%%%%%%%%%%%%%%%%%%%%%%%%%%%%%

\nocite{*}
\bibliographystyle{unsrt}
\bibliography{references}
\end{document}